\section{Discussion}
\label{sec:discussion}

% Author's note: this Section replaces the prior Discussion body. It drops the
% empirical-record recap, the "what the framework changes," and the regime
% subsections (those points are now made once, where they belong, in the
% Empirical results and Conclusion). The three subsections below earn their
% place: the measure family as the inferential object; economic significance;
% the general principle. "What skill means," "what the test does not claim,"
% and "conditioning is conservative" are each stated once and not repeated.

\subsection{The measure family as the inferential object}
\label{sec:discussion-measure}

The most serious objection to the framework is not that any particular verdict is
wrong but that the verdict is defined relative to a choice the analyst makes. The null
is indexed by a conditioning measure $\mathbb{Q}_\theta$ drawn from the family
$\mathfrak{Q}=\{\mathbb{Q}_\theta:\theta\in\Theta\}$ of Section~\ref{sec:theory}; all
members fix the realized path $\{P_t\}$ and structural profile $\mathcal{S}$, and they
differ only in which additional features of the realized schedule they hold fixed. That
choice is not pinned down by the data. It is a modeling decision on the same footing as
the specification of a regression or the choice of an instrument, and it should be
declared and defended as one. Each $\mathbb{Q}_\theta$ carries its own null $H_0^\theta$
and its own estimand $q_\theta=\mathbb{Q}_\theta\big(T(S)\ge T(S_0)\big)$, so what the
data deliver is not a single verdict but an \emph{identified set} of measure-specific
verdicts, one per admissible $\theta$. The object of inference is therefore that set, and
the analyst's task is to specify the admissible class $\Theta_0$ over which the set is to
be read---to declare, in advance, which partitions of the realized record into
``structure'' and ``placement'' are defensible for the question at hand.

Stating $\Theta_0$ in advance is what turns a free parameter into a model, and it is what
makes the partial-identification language of Section~\ref{sec:theory} operative rather
than decorative. The width of the identified set across $\Theta_0$ is the relevant
robustness statement: a verdict that is constant across the admissible class is sharply
identified, and one that flips is identified only as a set whose elements disagree. We
define \emph{measure-invariant entry-placement skill} as rejection of $H_0^\theta$ for
every $\theta\in\Theta_0$, and we adopt agreement across $\Theta_0$ as the bar a timing
claim must clear. The reason is not conservatism for its own sake but identification.
Because the measures differ precisely in where they draw the structure/placement line, a
rejection under $\mathbb{Q}_\theta$ but not $\mathbb{Q}_{\theta'}$ is confounded with the
conditioning choice $\theta$; only invariance to that choice isolates placement as
opposed to the analyst's partition. Measure invariance is thus a strictly stronger and
more credible claim than rejection under any single hand-picked measure, and it is the
right standard for a timing claim because the timing question is exactly the question that
the structure/placement partition is meant to answer.

Two boundaries of this standard must be stated plainly, and only once. Invariance is
\emph{necessary but not sufficient}: a confound shared by every admissible measure would
survive across $\Theta_0$, so agreement bounds the set of explanations but does not
collapse it to skill. And invariance is \emph{not necessary} for genuine ability, because
a conditioning measure can absorb real state-conditional skill into the structure it holds
fixed; a strategy that times entries well only within a market state which the measure
matches on will, correctly under that measure's null, fail to reject. The
schedule-measure sensitivity of Section~\ref{sec:schedule-sensitivity} is the canonical
illustration that these are features, not nuisances. The neutral gap-permutation measure
rejects none of the eleven strategies, while the context-matched measure---restricting
random entries to bars of comparable trailing volatility---rejects four of them
(Connors RSI2, ADX Trend, Breakout Vol+Mom, and Squeeze Breakout). The disagreement is
the substance. It locates the verdict's dependence on $\theta$ at a specific, interpretable
margin: whatever advantage those four strategies show is inseparable, on this evidence,
from a volatility-state alignment that the context-matched measure either credits as
placement or conditions away as structure, depending on which side of the line one draws.
Because no strategy rejects under both measures, the identified set of verdicts contains
both ``skill'' and ``no skill'' for those four, and the conjunctive standard returns the
honest reading: the absence of robust, measure-invariant entry-placement skill, with the
disagreement marking exactly where genuine state-conditional skill, if present, would have
to live. Measure-dependence is then not a defect of the framework to be apologized for but
the form in which the framework reports what the data can and cannot identify.

\subsection{Economic significance}
\label{sec:discussion-economic}

The finding carries a direct economic reading. Across the panel, entry-timing performance
is largely not separately identifiable from structural exposure to a drifting asset: to
first order the realized return is reproducible by randomized placement of the same trades.
For performance attribution and manager due diligence, a profitable track record is on its
own uninformative about timing skill, so an allocator should require a timing claim to
clear a structure-matched counterfactual before paying for it; the test supplies that
instrument, model-free and computable from a trade log. For the market-efficiency reading,
the result sharpens the object of study---rather than asking whether a strategy beats a
benchmark, it isolates the placement margin most often mistaken for skill and finds it
scarce and measure-fragile even where profitability is common. For strategy development, the
diagnostic redirects effort away from rules whose edge is conditioned away as structural
exposure, which will not survive live trading on a different path, and toward the regimes in
which placement evidence does concentrate (Section~\ref{sec:realistic-signals}), where
genuine if modest signal may reside.

\subsection{Beyond entry timing: a general principle}
\label{sec:discussion-general}

The construction is not specific to entry timing, to long-only trading, or to gold. It
applies whenever a decision sequence is overlaid on an exogenous process and one wishes to
separate the marginal contribution of the decisions from the structural exposure they
create: hold the structural footprint fixed, randomize only the decision margin, and refer
the realized statistic to the law this induces. Replacing ``entry placement on a price
path'' with the analogous decision--structure pair yields a test for factor-timing and
tactical asset-allocation rules (randomize the rebalancing schedule while holding portfolio
weights and the realized factor path fixed), for execution algorithms (randomize child-order
placement while holding the parent order and the realized microstructure path fixed), for
the evaluation of reinforcement-learning and agentic policies (randomize action timing while
holding the policy's structural footprint fixed, scoring against a structure-matched random
policy), and for the timing of macro forecasts and event signals (randomize the
announcement-relative placement). In each case the same three ingredients recur---a fixed
exogenous path, a preserved structural profile, and a randomized decision margin---so the
finite-sample validity of Proposition~\ref{prop:q-validity}, the local-power and
orthogonality results of Section~\ref{sec:theory}, and the calibration machinery transfer
without modification, and the measure family of Section~\ref{sec:discussion-measure} returns
as the object to specify. We develop the trading-timing instance in full and present these
as the program the framework opens.
